Возвращаемся к многомерной задаче. Шаг фиксирован: $\alpha_k = 1/L$ --- именно он максимизирует гарантированное убывание в \hyperref[note:379E4CAD-A44B-4631-8DAF-F3DC708EAEDE]{усиленном неравенстве}. Две теоремы ниже отличаются только классом функций, и таблица из определения невязки показывает, что меняется: без сильной выпуклости контролируется лишь норма градиента со скоростью $c/\sqrt{k}$, с ней --- сами точки, линейно со знаменателем $1 - \mu/L$.
